\honest{Belief in Intelligence}{Eliezer Yudkowsky}{2008}

I cannot predict Kasparov's moves. So what does my belief that he is a highly intelligent chess player predict? \nb{The post continues ``Expected Creative Surprises,'' posted the day before and not in the book.} To make it harder: if Kasparov plays a lesser grandmaster, Mr.~G, my guess at each man's move would be the same in any position, since all I can do is guess the best move with my own weak chess. \nb{If the probabilities for each next move are the same for both players in every position, they fix the chances of every whole game, and the results cannot favour Kasparov apart from the colour each plays. The day before, Yudkowsky had written that this situation is possible because ``I am not logically omniscient.'' This post leaves that out.}

My answer: the belief predicts the final position. It will be in the class of positions that are wins for Kasparov, and how strongly I believe he is better is how much probability I put on that class. The class is vague, but more specific than total ignorance, since a vast set of final positions would falsify it. I know where Kasparov is trying to steer the game, and I expect he is strong enough to get there, though not how. \nb{Chess ratings work this way: under the Elo system a rating difference predicts an expected score and no moves.}

Or take a friend who drives me to the airport in a city I do not know. I cannot predict a single turn, yet I can predict that we will arrive, from wherever we start, before I get into the car; with a flight to catch, I would not get in otherwise. ``Isn't this a remarkable situation to be in, from a scientific perspective?'' \nb{William James had described it in 1890: for iron filings drawn to a magnet the path is fixed, but for Romeo seeking Juliet ``it is the end which is fixed, the path may be modified indefinitely.'' Daniel Dennett's ``intentional stance'' makes the same kind of prediction from an agent's goals. The post mentions neither.}

Usually we predict by running a model forward step by step, as for the planets, or by a closed-form solution, as for a coin that turns over every minute, where the same formula gives the face after a hundred minutes and at every minute between. My airport model predicts the end, predicts no step, and needs no starting point. It needs only my friend's goal, planning skill and knowledge of the city, the last a very abstract relation between friend and city that needs no specific knowledge of either.

Understanding this strange abstract knowledge is, in a sense, my life's work. ``Intelligence'' is too narrow a word for these situations; I would say ``optimization process.'' Natural selection is one, since ``we can't predict the exact form of the next organism observed.'' \nb{That gives the half of the example that cannot be predicted, and not the outcome that can.} My specialty is ``Friendly Artificial Intelligence,'' of which, I hope, I will gain especially precise abstract knowledge.
